\chapter{Introduction générale}

\section{Contexte du projet}
Le présent travail s'inscrit dans un projet de fin d'études consacré à l'aide à
la décision sur le marché actions marocain. Dans un tel contexte, l'utilisateur
ne recherche pas uniquement une visualisation de marché, mais un dispositif
capable de produire des signaux techniques cohérents, traçables et directement
mobilisables dans une lecture quotidienne des titres suivis.

Le besoin fonctionnel ayant motivé cette partie du projet peut être formulé de
manière simple : disposer d'une chaîne fiable permettant de transformer des
données de marché normalisées en indications de lecture synthétiques, tout en
conservant la possibilité d'expliquer l'origine des résultats affichés. Ce
double objectif, à la fois quantitatif et ergonomique, justifie l'attention
accordée à la robustesse méthodologique ainsi qu'à la qualité de restitution
dans l'interface.

\section{Problématique}
La difficulté principale ne réside pas dans la production d'un indicateur
isolé, mais dans la conception d'un mécanisme de sélection suffisamment solide
pour limiter le sur-ajustement et rendre les résultats lisibles pour un usage
opérationnel. Dès lors, la problématique traitée dans ce document peut être
formulée comme suit :

\begin{quote}
\textit{Comment construire des signaux techniques robustes sur un marché bruité,
puis les restituer dans une interface de manière suffisamment claire pour
soutenir une prise de décision quotidienne ?}
\end{quote}

Cette problématique se décompose en trois enjeux. Le premier concerne le biais
de sélection : lorsqu'un grand nombre de variantes est testé, une partie des
performances observées peut relever du hasard plutôt que d'une structure
réellement exploitable \cite{deprado2018,bailey2014}. Le deuxième tient à
l'instabilité temporelle : un signal performant sur une période donnée peut se
détériorer dès que les conditions de marché changent \cite{pardo2008}. Le
troisième relève de l'explicabilité : une restitution purement agrégée perd une
part importante de sa valeur si l'utilisateur ne peut pas comprendre quels
éléments la composent.

\section{Objectifs}
Cette version du mémoire poursuit quatre objectifs complémentaires :
\begin{enumerate}
\item définir une méthodologie de génération et de filtrage des signaux qui
privilégie la validation hors échantillon ;
\item formaliser la manière dont plusieurs familles d'indicateurs sont
transformées en résultats comparables et synthétiques ;
\item décrire la version initiale de la page Signaux comme espace d'analyse
progressive, allant d'une vue d'ensemble vers le détail ;
\item montrer comment les résultats sont ensuite restitués dans le tableau de
bord afin de faciliter la hiérarchisation des opportunités.
\end{enumerate}

\section{Démarche retenue}
Le rapport suit une logique unidirectionnelle :
\[
\text{Données OHLCV validées} \rightarrow \text{Méthodologie de signal}
\rightarrow \text{Page Signaux} \rightarrow \text{Tableau de bord}
\]

Cette organisation permet de séparer clairement la production du résultat, son
interprétation analytique et sa restitution décisionnelle. Elle évite
également de confondre la qualité d'un affichage avec la validité de la méthode
ayant conduit au signal.

\section{Organisation du document}
Le deuxième chapitre expose le cadre conceptuel qui justifie les choix de
conception. Le troisième présente la méthodologie de détection des signaux. Le
quatrième décrit la version initiale de la page Signaux. Le cinquième explique
la représentation des résultats dans le tableau de bord. Enfin, le sixième
chapitre propose une conclusion partielle sur l'apport de cette chaîne dans le
projet global.

